\documentclass[11pt,a4paper]{article}
\usepackage[T1]{fontenc}
\usepackage[utf8]{inputenc}
\usepackage[brazilian]{babel}
\usepackage{lmodern}
\usepackage[margin=2.3cm]{geometry}
\usepackage{enumitem}
\usepackage{xcolor}
\usepackage{fancyhdr}

\pagestyle{fancy}
\fancyhf{}
\rhead{\small Pesquisa acadêmica, UNISINOS}
\cfoot{\small\thepage}
\renewcommand{\headrulewidth}{0pt}

\setlength{\parindent}{0pt}
\setlength{\parskip}{0.55em}
\setlist[enumerate]{leftmargin=*, itemsep=0.25em, topsep=0.2em}
\setlist[itemize]{leftmargin=*, itemsep=0.2em, topsep=0.2em}
\newcommand{\tema}[1]{\subsection*{\normalsize #1}}

\begin{document}

\begin{center}
  {\Large\bfseries Cardápios digitais em bares e restaurantes de Porto Alegre}\\[0.3em]
  {\large Convite para entrevista com fornecedor de tecnologia}
\end{center}

\section*{\normalsize Sobre a pesquisa}

Sou César Hoffmann, aluno de Ciência da Computação na UNISINOS, e estou concluindo meu Trabalho de Conclusão de Curso, orientado pelo Prof.\ Dr.\ Guilherme Lacerda. A pesquisa investiga como bares e restaurantes de Porto Alegre adotaram, ou escolheram não adotar, cardápios digitais e autoatendimento desde 2020, e como os clientes vivenciam essa mudança.

Até agora, ouvimos gestores de estabelecimentos com perfis diferentes e mais de cem frequentadores de bares e restaurantes. Falta a visão de quem constrói e mantém essas soluções: como elas são pensadas, o que os estabelecimentos pedem, o que funciona bem e o que ainda é difícil. É essa perspectiva que gostaríamos de ouvir de vocês.

\section*{\normalsize Como seria a conversa}

\begin{itemize}
  \item Cerca de 30 minutos, por videoconferência, no dia e horário que for melhor para você.
  \item Gravada apenas para transcrição, com o seu consentimento. A gravação não é divulgada.
  \item O nome da empresa só aparece no trabalho se você autorizar; caso contrário, usamos um pseudônimo. Seu nome pessoal não é divulgado.
  \item Não é preciso preparar nada nem compartilhar números sigilosos. Qualquer pergunta pode ser pulada.
  \item Se for útil, assinamos também um termo de confidencialidade da UNISINOS.
  \item Ao final, envio um resumo dos resultados da pesquisa, que pode ser útil para o time de produto.
\end{itemize}

\section*{\normalsize Temas e perguntas}

As perguntas abaixo servem de guia. Não precisamos passar por todas em 30 minutos: a conversa é livre, e os temas 2 e 4 são os mais importantes para a pesquisa.

\tema{1. A empresa e o produto}
\begin{enumerate}
  \item Como surgiu a solução e que problema ela queria resolver?
  \item Qual é o seu papel na empresa?
  \item Que tipo de estabelecimento usa a solução, em Porto Alegre e no Rio Grande do Sul?
  \item Que modalidades de cardápio digital e autoatendimento vocês oferecem (QR Code só para consulta, QR Code com pedido, tablet, totem)? Quais são as mais usadas e por quê?
\end{enumerate}

\tema{2. Como a solução funciona}
\begin{enumerate}[resume]
  \item Em linhas gerais, como a solução é construída e hospedada para atender muitos estabelecimentos ao mesmo tempo?
  \item Com que outros sistemas ela se integra (caixa, nota fiscal, estoque, cozinha, pagamento, delivery)? Quais integrações dão mais trabalho?
  \item O que acontece quando a internet do estabelecimento ou do cliente cai?
  \item Como o estabelecimento atualiza preços e itens esgotados, e quanto tempo isso leva para aparecer ao cliente?
  \item O que cada estabelecimento pode personalizar e o que é igual para todos?
\end{enumerate}

\tema{3. Necessidades dos estabelecimentos e dos clientes}
\begin{enumerate}[resume]
  \item Quais funcionalidades os estabelecimentos mais pedem? E quais quase não usam?
  \item Que retorno vocês recebem de quem está na mesa? Quais são as reclamações e os elogios mais comuns?
  \item Alguma funcionalidade não deu certo? O que vocês aprenderam com isso?
\end{enumerate}

\tema{4. Dados pessoais e acessibilidade}
\begin{enumerate}[resume]
  \item O cliente final precisa informar algum dado pessoal para ver o cardápio ou pedir? Quem decide isso?
  \item Como vocês tratam a LGPD na relação com o estabelecimento?
  \item Como a solução lida com acessibilidade (tamanho de fonte, contraste, leitores de tela, pessoas idosas)?
\end{enumerate}

\tema{5. Mercado}
\begin{enumerate}[resume]
  \item Que tipo de estabelecimento adota com mais facilidade e qual resiste? Quais são as objeções mais comuns?
  \item O que mudou na procura pela solução durante e depois da pandemia?
  \item A lei de Porto Alegre de 2024, que exige cardápio físico junto com o digital, teve algum efeito para vocês?
  \item Daqui a cinco anos, como você imagina o atendimento em bares e restaurantes?
\end{enumerate}

\vspace{0.6em}
\textbf{Contato:} César Hoffmann \textbullet{} cesarhoff@unisinos.br \textbullet{} (51) 99804-4896 (WhatsApp)\\
\textbf{Orientador:} Prof.\ Dr.\ Guilherme Lacerda, Escola Politécnica, UNISINOS

\end{document}
